\begin{helper}
Fix \(n,m,k\) and \(\varepsilon\), and assume \(0<n\).  Define
\[
S=\left\lceil
\frac{k n^{1/4-\varepsilon}+\noderef{TwoBiteLargeCutoff}(\varepsilon,n)}
     {\noderef{TwoBiteSmallCutoff}(\varepsilon,n)}
\right\rceil .
\]
Then every finite witness family \(B\) of two-bite base vertices satisfying
\[
|B|\,\noderef{TwoBiteSmallCutoff}(\varepsilon,n)
\le
k n^{1/4-\varepsilon}+\noderef{TwoBiteLargeCutoff}(\varepsilon,n)
\]
has \(|B|\le S\).  Moreover
\[
S\le k n^{1/4-3\varepsilon}+n^{1/4-\varepsilon}+1 .
\]
\end{helper}
\begin{proof}
Since \(0<n\), the small cutoff
\(\noderef{TwoBiteSmallCutoff}(\varepsilon,n)=n^{2\varepsilon}\) is positive.
Dividing the displayed witness-size inequality by this positive quantity gives
\[
|B|\le
\frac{k n^{1/4-\varepsilon}+\noderef{TwoBiteLargeCutoff}(\varepsilon,n)}
     {\noderef{TwoBiteSmallCutoff}(\varepsilon,n)}.
\]
The definition of \(S\) as the ceiling of the right hand side therefore gives
\(|B|\le S\).

Using the cutoff definitions,
\[
\noderef{TwoBiteLargeCutoff}(\varepsilon,n)=n^{1/4+\varepsilon},
\qquad
\noderef{TwoBiteSmallCutoff}(\varepsilon,n)=n^{2\varepsilon}.
\]
The same positivity assumption permits the usual exponent subtraction laws, so
\[
\frac{k n^{1/4-\varepsilon}+n^{1/4+\varepsilon}}{n^{2\varepsilon}}
= k n^{1/4-3\varepsilon}+n^{1/4-\varepsilon}.
\]
Finally, the ceiling of a nonnegative real number is less than that number plus
one, hence at most that number plus one.  This gives the stated upper bound on
\(S\).
\end{proof}
